\documentclass[11pt,a4paper]{article}

\usepackage[utf8]{inputenc}
\usepackage[T1]{fontenc}
\usepackage[margin=2.5cm]{geometry}
\usepackage{amsmath}
\usepackage{booktabs}
\usepackage{graphicx}
\usepackage{siunitx}
\usepackage{hyperref}

\title{Detection and Classification of Aspirated Tissue Pieces\\
       in Marginum Fluorescence Recordings}
\author{IDP2026 Group 8}
\date{\today}

\begin{document}

\maketitle

\begin{abstract}
% TODO: summarise the problem, the method and the main results in 150--200 words.
We present a pipeline that reads large two-channel Marginum fluorescence recordings
(\SI{50}{\kilo\hertz}, up to several gigabytes), removes noise and baseline drift,
detects aspirated tissue pieces and labels each one as dyed tissue, undyed tissue or water.
An interactive Streamlit application lets the user browse the recording and inspect every
detected piece.
\end{abstract}

\tableofcontents

%----------------------------------------------------------------------
\section{Introduction}
% TODO: background on Marginum, why detecting tissue pieces matters, project goals.

The recordings are CSV files with two integer columns, \texttt{adc1} and \texttt{adc2},
sampled at \SI{50}{\kilo\hertz}. The goal is to find each aspirated tissue piece,
decide whether it is tissue (with or without dye) or water, and estimate the amount of tissue.

%----------------------------------------------------------------------
\section{Data}
% TODO: describe the measurement sessions (2023 test signals, 2025 known-volume and tumour recordings).

%----------------------------------------------------------------------
\section{Methods}

\subsection{Fast reading and reduction to \SI{1}{\kilo\hertz}}
The file is split into 1024 blocks whose lines are counted in parallel to build a row index,
so any time window can be read directly. The blocks are then reduced to the mean, minimum and
maximum of every 50 samples (\SI{50}{\kilo\hertz} $\rightarrow$ \SI{1}{\kilo\hertz}).
The signal has no content above about \SI{500}{\hertz}, so the \SI{1}{\milli\second} mean keeps
99.3--99.9\,\% of the peak height. The index and summary are cached for later opens.

\subsection{Preprocessing (Part A)}
For each channel:
\begin{enumerate}
  \item median filter (3 samples) to remove single-sample spikes;
  \item Savitzky--Golay filter (7 samples, order 3);
  \item subtraction of a baseline fitted by asymmetric least squares (ALS) on \SI{10}{\hertz}
        block means and interpolated back to \SI{1}{\kilo\hertz};
  \item noise level $\sigma$: RMS of the standard deviation of the quietest 10\,\% of
        \SI{200}{\milli\second} blocks, after dropping idle blocks.
\end{enumerate}
If the median ratio \texttt{adc1}/\texttt{adc2} at the 50 largest \texttt{adc2} peaks is below
0.1, \texttt{adc1} is not used for the two-channel check.

\subsection{Detection and classification (Part B)}
Peaks with prominence above $10\sigma$ are detected on both channels. Peaks less than
\SI{1}{\second} apart are grouped into one piece. Each piece is labelled from its largest peak
(Table~\ref{tab:labels}).

\begin{table}[h]
  \centering
  \caption{Labelling rules.}
  \label{tab:labels}
  \begin{tabular}{ll}
    \toprule
    Condition & Label \\
    \midrule
    above the no-dye level (68.9 counts) & \texttt{dyed} \\
    below, confirmed by \texttt{adc1} & \texttt{solid\_no\_dye} \\
    below, not confirmed & \texttt{water} \\
    below, two-channel check not usable & \texttt{undecided} \\
    \bottomrule
  \end{tabular}
\end{table}

The two-channel check requires an \texttt{adc1} peak within \SI{\pm 20}{\milli\second} of the
top \texttt{adc2} peak (shifted by the measured channel lag) and at least 5\,\% as high.
The chance rate is estimated by shifting \texttt{adc1} by 2--\SI{20}{\second}; above 20\,\%
the check is disabled.

\subsection{Amount of tissue}
The amount is the total area (counts\,$\times$\,ms) of the \texttt{dyed} and
\texttt{solid\_no\_dye} pieces. No conversion to millilitres is made yet, since the area per
millilitre depends on the dye concentration and on re-suction.

\subsection{Visualizer}
% TODO: describe the Streamlit app (overview with draggable window, detail view, heatmap, events table, LTTB).

%----------------------------------------------------------------------
\section{Results}

\subsection{Speed}
\begin{table}[h]
  \centering
  \caption{Processing time for a \SI{7.6}{\giga\byte}, \SI{4.7}{\hour} recording (MacBook, 4 workers).}
  \label{tab:speed}
  \begin{tabular}{lrr}
    \toprule
    Step & First open (s) & Next open (s) \\
    \midrule
    CSV $\rightarrow$ \SI{1}{\kilo\hertz} & 7.5 & 0.1 \\
    Part A & 1.8 & 1.8 \\
    Part B & 1.8 & 1.8 \\
    \midrule
    Total & 11 & 3.7 \\
    \bottomrule
  \end{tabular}
\end{table}

\subsection{Detection results}
% TODO: add the per-recording table from result.json and discuss.

%----------------------------------------------------------------------
\section{Discussion}
% TODO: limitations (merged pieces in dense suctions, ml calibration, undecided recordings), future work.

%----------------------------------------------------------------------
\section{Conclusion}
% TODO

%----------------------------------------------------------------------
\section*{Contributions}
% TODO: who did what (see README.md).

\end{document}
